\subsection*{分区身份与当前承载}

第十六章把会话、受理绑定、付款恢复和消息发布放在购买权威一侧。现在把这份权威拆成三个逻辑分区 b1、b2、b3。示例用账户编号确定归属：账户 1、4、7 进入 b1，账户 2、5、8 进入 b2，账户 3、6、9 进入 b3。分区函数依赖业务身份中稳定的账户编号，客户端重试时沿用原账户和完整 $S$，便会落到相同的逻辑分区。这里的完整身份继续是
\[
S=(\texttt{account\_id},\texttt{action},
   \texttt{session\_id},\texttt{idempotency\_key}) .
\]
账户编号决定存放位置，动作、会话和键共同参与原操作的识别。把其中一个分量丢掉，得到的就已经是另一个仲裁范围。

b1、b2、b3 描述业务事实的归属，G1、G2、G3 描述当前保存事实的副本组。初始目录把 b1 指向 G1，后续可以把 b1 搬到 G2，b1 的编号和 Alice 保存的 $S$ 都保持原值。副本组还可以在内部替换一台服务器，此时目录中的 $\mathrm{b1}\to\mathrm{G1}$ 可以继续成立。组内换副本改变复制成员，业务分区换组改变当前承载，管理者根据这两个不同对象选择操作。

\CSeventeenFigure{ch17-01-ownership}{逻辑分区保存稳定归属，目录记录副本组的当前承载}{fig:ch17-ownership}

入口从目录读到 \code{\{owner:G1, epoch:1\}} 后，可以缓存这条路由。请求带着 b1 和归属代次 1 到达 G1，由 G1 内部的 Raft leader 处理。目录帮助入口找到地址，G1 对分区模式和归属代次的核验决定本次修改是否有资格进入业务状态机。正常请求在缓存有效时直接访问所属组，控制面负责配置变化，无须为每一次购买分配一把全局锁。

如果入口缓存过期，它可能继续把请求发给旧组。旧组的答复提供重新定位的线索，入口刷新目录，再把原 $S$ 和原输入送到当前组。这条恢复路径中变化的是地址和归属代次。把 \code{idempotency\_key} 换成新值，会把一次位置变化扩大成另一项业务操作。

\subsection*{业务分区是一个闭合的数据集合}

b1 中已受理的 $S$ 绑定到原输入、\code{P:b1:1}、\code{C:b1:1} 和 \code{T:b1:1}。第十四至十六章中承担权威受理职责的 $A_J$，在这里属于该逻辑分区的权威状态。支付工作者读取这一绑定，按原 \code{provider\_scope} 和 \code{pay-C:b1:1} 调用平台。平台返回 X1 后，分区保存完整付款证据、原付款结果、订单 \code{O:b1:1} 和事件 \code{EV:b1:1}，由 $R_J$ 固定这些关联。事件原报文属于待发布事实，消息平台里的消息副本属于后续传递，迁移时按各保存域交接责任。

分区还保存 Alice 的会话期限与关闭状态、工作者当前代次、尚未解决的付款待办、事件发布标记、消费交接回执与当前履约进度，以及退款所引用的原扣款和占额。若购买包含跨分区库存预留，购买分区保存协调者记录，库存所属分区保存参与者记录。表~\ref{tab:ch17-state} 列出购买分区自身的闭合范围。

\begin{table}[htbp]\centering\small
\begin{tabularx}{\linewidth}{@{}p{.19\linewidth}XX@{}}
\toprule
分区中的对象 & 保存的联系 & 后续读取与修改\\
\midrule
会话与受理绑定 & 会话准入边界、完整 $S$、原输入版本、P/C/T、原平台范围与键 & 入口核验准入，已有绑定比较输入，服务端恢复原受理\\
付款结果与待办 & 平台原输入、执行代次、原证据、原结果、O/EV、发布标记、交接回执及来源 & 工作者接管、证据归并、查询恢复、Outbox 发布与交接停扫\\
退款与事务责任 & 原 C/X、退款身份与占额、库存状态、prepare 记录、协调者决定 & 退款恢复、容量判断、库存完成与迁移准入\\
\bottomrule
\end{tabularx}
\caption{购买逻辑分区保存的完整业务状态}\label{tab:ch17-state}
\end{table}

这些对象之间有可检查的联系。退款记录引用的 C 必须能在同一业务历史中找到原扣款，已发布标记必须对应保存下来的原事件，某个 \code{pending} 付款必须带有可以重发或查询的原平台范围与键。只复制一张 $S\to$结果的表，会丢失尚未取得结果的操作；只复制当前订单，又会丢失会话准入和退款占额。闭合范围由系统已经承担的恢复责任确定。

\CSeventeenFigure{ch17-02-closure}{分区图像连接身份、原结果与尚未完成的业务责任}{fig:ch17-closure}

原受理、原结果和当前进度也分别保存。$A_J$ 已固定 C，后来换了几次工作者，这份受理继续指向原 C；$R_J$ 已固定成功结果，后来完成履约，原付款结果保持原值。当前进度可以从待付款前进到已付款、已履约，查询把它作为另一项观察返回。快照和迁移同时带走固定关联与可变进度，新位置才有足够事实回答不同用途的读取。

\CSeventeenFigure{ch17-15-facts-progress}{原受理与原结果保持原关联，当前进度沿后续历史前进}{fig:ch17-facts-progress}

本书为新操作采用稳定分区编号与分区内序号组成的编号。b1 的序号从 1 前进到 2，得到 \code{C:b1:2}；b1 从 G1 搬到 G2 后，新的受理继续使用 b1 的序号空间。序号增加和受理绑定在同一次状态机应用中保存，客户端看到已提交受理时，相应编号已经属于该操作。

旧系统导入的 C1、P1、T1 和 \code{pay-C1} 按原关联保存。导入器可以为记录增加承载位置和格式版本，原平台键继续指向原平台操作。若旧记录已预留 O1、EV1，后续补全付款结果时也沿用这两个编号。根据新组名重新生成 C，或根据新编号格式猜测旧订单，会破坏已经存在的外部关联。

\CSeventeenFigure{ch17-16-id-continuity}{新编号沿稳定分区分配，旧编号随原关系导入}{fig:ch17-id-continuity}

序号有一个来源条件：恢复出来的业务历史包含此前已经对外承诺的分配。整个副本组故障后若从旧备份重建，旧序号和旧绑定可能一并消失。第十四章由不可回退来源提供灾难代次，并先冻结、核对外部证据，取回原 $S\to$C1$\to$\code{pay-C1} 关联，随后才按开放范围接受新意图。把备份计数器简单加一，无法判断这个值是否曾在丢失的后缀中发出。正常复制与灾难回退因此使用不同的恢复依据。
\FloatBarrier
